\section{Dykstra's projection algorithm} \label{sec:dykstra_algorithm}

The MAP algorithm is an iterative method to find a point $x_m$ in the intersection region, $x_m  \in\bigcap_{i=0}^{n-1}\mathcal{H}_i$, by cyclically projecting onto the convex sets $\mathcal{H}_i$~\cite{NEUMANN,BREGMAN1965}. Starting from an arbitrary point $x_0=x^\circ$, MAP generates a sequence of iterates $\{x_m\}$ by computing successive projections onto each of the half-spaces,
\begin{align}\label{eq:map}
    x_{m+1} = (P_{\mathcal{H}_{n-1}} \circ \dots \circ P_{\mathcal{H}_0}) (x_m).
\end{align}
When each $\mathcal{H}_i$ is an affine subspace, it can be shown that the MAP also solves problem~\eqref{eq:projection}~\cite{NEUMANN}. For general closed convex sets, however, the sequence generated by successive iterations of~\eqref{eq:map} is only guaranteed to converge to a point within the intersection $\mathcal{H}$, and not to the Euclidean projection $x^\star$~\cite{BREGMAN1965}.

Dykstra's method extends~\eqref{eq:map} by introducing auxiliary variables $e_m \in \mathcal{X}$ that store the negative error vectors associated with the projection of a primary variable $x_m$ onto a given half-space. The algorithm proceeds by generating a series of primary iterates \{$x_{m}$\} and auxiliary iterates \{$e_{m}$\} using the following scheme:
\begin{subequations}\label{eq:dykstra}
\begin{align}
x_{m+1}&=\mathcal{P}_{\mathcal{H}_{[m]}}\left(x_{m}+e_{m-n}\right),\label{eq:dykstra:proj}\\
e_m&=e_{m-n}+x_{m}-x_{m+1}\label{eq:dykstra:error},
\end{align}
\end{subequations}
where $[m]$ represents the modulus operator $[m]\eqdef m \, \text{mod} \, n$. The decision variable is initialised as $x_0=x^\circ$ and the auxiliary variables $e_m$ as $e_{-n} = e_{-(n-1)} = \dots = e_{-1} = 0$.
The Boyle-Dykstra theorem~\cite{DYKSTRA,DYKSTRAPERKINS} implies asymptotic convergence to the projection, so that $\lim_{m\rightarrow\infty}\anynorm{x_m-\mathcal{P}_\mathcal{H}(x^\circ)}=0$. For any finite number of iterations $m$, there is no guarantee that the iterates $x_m$ lie in the intersection $\mathcal{H}$, nor that $x_m\neq x^\circ$. Note that the MAP update~\eqref{eq:map} can be recovered from Dykstra's update~\eqref{eq:dykstra} by setting the auxiliary variables to zero in~\eqref{eq:dykstra:error}, i.e.\ $e_m=0$ for all $m$.

For polyhedral sets as defined in~\eqref{eq:sets_i}, the projection step~\eqref{eq:dykstra:proj} of Dykstra's method can be simplified to
\begin{align}\label{eq:dykstra:proj:poly}
x_{m+1}=
\begin{cases}
x_{m}+e_{m-n} & \text{if } x_{m}+e_{m-n}\in\mathcal{H}_{[m]}\\
x_{m} - \left(x_{m}^T a_{[m]} - b_{[m]}\right) a_{[m]} & \text{otherwise}
\end{cases},
\end{align}
and the update for the auxiliary vector~\eqref{eq:dykstra:error} to
\begin{align}\label{eq:dykstra:error:poly}
e_m=
\begin{cases}
0 & \text{if } x_{m}+e_{m-n}\in\mathcal{H}_{[m]}\\
e_{m-n}+\left(x_{m}^T a_{[m]} - b_{[m]}\right) a_{[m]} & \text{otherwise}
\end{cases}.
\end{align}
The auxiliary vector $e_m$ is either zero or parallel to the half-space norm $a_{[m]}$, so it can be characterised via a scalar quantity $k_m \in \mathbb{R}$ as $e_m = k_m a_{[m]}$ with 
\begin{align}\label{eq:km}
k_m = 
\begin{cases}
0 & \text{if } x_{m}+k_{m-n}a_{[m]}\in\mathcal{H}_{[m]}\\
k_{m-n} + x_m^T a_{[m]} - b_{[m]} & \text{otherwise}
\end{cases}.
\end{align}
The convergence of Dykstra's iterates to the Eucledian projection has been analysed in~\cite{DYKSTRAPERKINS,DYKSTRAPOLY,DYKSTRAPOLY2} for polyhedral sets. The analysis in~\cite{DYKSTRAPOLY} is based on partitioning the collection $\{\mathcal{H}_i\}_{i=0}^{n-1}$ into an \emph{inactive} subset---half-spaces containing the projection $x^\star$ in their interior $\text{int}\,\mathcal{H}_i$---and an \emph{active} subset---half-spaces containing the projection on their boundary $\partial\mathcal{H}_i$. The two subsets, $\mathcal{A}$ and $\mathcal{B}$, respectively, are defined as:
\begin{align}
&\mathcal{A} \eqdef \set{i\in\lbrace 0,\dots,n-1\rbrace}{x_\infty\in \partial\mathcal{H}_i},
&\mathcal{B}\eqdef \lbrace 0,\dots,n-1\rbrace\backslash \mathcal{A}=\set{i\in\lbrace 0,\dots,n-1\rbrace}{x_\infty\in \text{int}\,\mathcal{H}_i},
\end{align}
where $x_\infty=\lim_{m\rightarrow\infty}x_m$. It can be shown that there exists a number $N_1$ such that when $[m]\in \mathcal{B}$ at iteration $m\geq N_1$, it follows that $x_m=x_{m-1}$ and $e_m=0$, i.e.\ the half-spaces that become inactive remain so~\cite[Lemma~3.1]{DYKSTRAPOLY}. Furthermore, there exists $N_2\geq N_1$ such that whenever $n\geq N_2$, it holds that $\twonorm{x_{m+n}-x_\infty}\leq\alpha_{[m]}\twonorm{x_m-x_\infty}$, where $0\leq\alpha_{[m]}<1$ are scalars related to angles between half-spaces~\cite[Lemma~3.7]{DYKSTRAPOLY}. The integer $N_2$ describes the iteration from which on the algorithm has determined the inactive half-spaces. It then follows that there exist constants $0\leq c < 1$ and $\rho > 0$ such that $\twonorm{x_m -x_\infty} \leq \rho\, c^m$~\cite[Thm.~3.8]{DYKSTRAPOLY}. The constant $c$ can be estimated from the smallest $\alpha_{[m]}$, which is characterised by the angle between certain subspaces (subspaces formed by active half-spaces). The constant $\rho$, however, depends on an unknown $N_3\geq N_2$ and on $x^\circ$, and can therefore not be computed in advance~\cite{DYKSTRAPERKINS,XUPOLY}.
% In the case of stalling, the variable $\rho$ can become arbitrarily large, making the application of Dykstra's method difficult in practice. The authors of~\cite{DYKSTRAPERKINS} proposed a combined Dysktra-conjugate-gradient method that allows for computing an upper bound on $\anynorm{x_m -x_\infty}$. The authors of~\cite{XUPOLY} proposed an alternative algorithm called \emph{successive approximate algorithm}, which promises fast convergence, conditioned on knowing a point $x\in\mathcal{H}$ in advance.

% \subsection{Fenchel Dual and ADMM} \label{subsec:admm_connection}

% Since its original formulation, dykstra's method has been connected with several areas of optimisation in recent years. In particular, it is possible to reformulate the Hilbert space BAP using the sum of the indicator functions $\delta_{\mathcal{H}_i}(x)$ for each constraint set, yielding the augmented unconstrained optimisation problem
% \begin{equation} \label{eq:augmented_bap}
%     \underset{x \in \mathbb{R}^p}{\text{minimise}} \quad  \frac{1}{2}\|x - x^\circ\|^2 + \sum_{i=0}^{n-1} \delta_{\mathcal{H}_i}(x),
% \end{equation}
% \[
% \delta_{\mathcal{H}_i}(x) \eqdef
% \begin{cases}
% 0 & \text{if } x \in \mathcal{H}_i, \\
% +\infty & \text{if } x \notin \mathcal{H}_i,
% \end{cases}
% \]
% The sequence of iterates generated by Dykstra's algorithm on primal problem~\eqref{eq:augmented_bap} was shown to correspond to the sequence generated by a cyclic block coordinate descent algorithm on its Fenchel dual~\cite{HAN1988,GAFFKE1989}.
% This equivalence reveals that the auxiliary increment sequences $\{e_m\}$ in the primal formulation are precisely the dual variables of the optimisation problem. Consequently, the update rule~\eqref{eq:dykstra:error} can be interpreted as the execution of a coordinate-wise update step in the dual space. From this dual perspective, it is possible to transfer analytical results between the two formulations: known convergence rates for Dykstra's algorithm when applied to polyhedral sets can be used to establish linear convergence for coordinate descent on dual problems, such as the LASSO in machine learning.
% The connection to duality also situates Dykstra's algorithm within the broader family of operator splitting methods in close relationship with the Alternating Direction Method of Multipliers (ADMM)~\cite{BOYDADMM}. The ADMM algorithm decomposes the problem via an augmented Lagrangian, which adds a quadratic penalty on constraint violations to the standard Lagrangian (as opposed to infinite costs with the indicator function reformulation). The algorithm then proceeds by performing alternating minimisation steps with respect to the primal variables, followed by an update of the dual variables. This structure differs from the purely dual ascent nature of Dykstra's algorithm. However, for the specific case of the Euclidean BAP with two constraint sets ($\mathcal{H}_0, \mathcal{H}_1$), it has been shown that if one of the sets is a linear subspace, then applying ADMM to the dual problem generates iterates that are equivalent to those produced by Dykstra's algorithm on the primal problem~\cite{GABAY1983}. In more general cases, however, the algorithms are distinct. Dykstra's method performs a cyclic, coordinate-wise update in the dual space, whereas ADMM's updates are typically not coordinate-wise and its convergence behaviour is sensitive to the choice of the penalty parameter.

% \subsection{The Stalling Phenomenon} \label{subsec:stalling}

Although the Boyle-Dykstra theorem implies guaranteed asymptotic convergence to the optimal solution $x^\star$, the iterates in Dykstra's algorithm are not guaranteed to converge monotonically. It is in fact possible to establish initial conditions for which the iterates exhibit a cyclic repetition pattern, known as \emph{stalling}.
As a motivating example, in~\cite{DYKSTRASTALLING} the behaviour of Dykstra's method is analysed for two polyhedral sets in $\mathbb{R}^2$, namely an origin-centred box intersected with a line. For this specific example, the authors give conditions on Dykstra's algorithm for (i) finite convergence, (ii) infinite convergence, and (iii) stalling followed by infinite convergence, identifying sets of initial conditions to achieve each of the convergence mechanisms.